\section{Retroactive Public Goods Funding}

\label{sec:retro}
% ══════════════════════════════════════════════════════════════════════════════

\subsection{Motivation}

Prospective funding rewards proposals; retroactive funding rewards
results. This reduces the risk of funding bad proposals and increases
the reward for genuine public goods \citep{boudreau2016looking}.

\subsection{Mechanism}

Quarterly, the DAO allocates $R_{\text{pool}}$ (20\% of budget) for
retroactive funding. Eligible outputs: open-source software used by
$\geq$5 projects, datasets with $\geq$100 uses, papers with $\geq$20
citations or verified real-world impact, and infrastructure with
$\geq$6 months uptime.

Allocation uses QV with a 21-member badge holder committee ($\geq$1000
ZooRep), following the Optimism retroPGF model
\citep{optimism2023retropgf} adapted for scientific impact.

\subsection{Impact Metrics}

\begin{equation}
  I = w_1 \cdot \text{Reuse} + w_2 \cdot \text{Citation} +
  w_3 \cdot \text{Conservation} + w_4 \cdot \text{Community}
\end{equation}

Weights set by the Research Council and updated through ZIPs.


% ══════════════════════════════════════════════════════════════════════════════
